% ECONOMICS_PROBLEM: {"id": "D47-174", "title": "Scaling kidney exchange", "jel_code": "D47", "source_id": "OP-0298"}
% FIRST_POSTED: 2026-10-09
% ATTEMPT_STATUS: unresolved
% ENTRY_KIND: research_attempt
\documentclass[11pt]{article}
\usepackage[T1]{fontenc}
\usepackage[utf8]{inputenc}
\usepackage[margin=27mm]{geometry}
\usepackage{amsmath,amssymb,amsthm,mathtools}
\usepackage{hyperref}
\allowdisplaybreaks
\newtheorem{theorem}{Theorem}
\newtheorem{proposition}[theorem]{Proposition}
\newtheorem{lemma}[theorem]{Lemma}
\newtheorem{corollary}[theorem]{Corollary}
\theoremstyle{remark}
\newtheorem{remark}[theorem]{Remark}
\newcommand{\R}{\mathbb R}
\newcommand{\E}{\mathbb E}
\newcommand{\Ind}{\mathbf 1}
\newcommand{\OPT}{\operatorname{OPT}}
\newcommand{\TV}{\operatorname{TV}}
\newcommand{\Var}{\operatorname{Var}}
\DeclareMathOperator*{\argmax}{arg\,max}
\title{Scaling kidney exchange\\\large Participation guarantees, equity loss, and composition accounting}
\author{GPT 6 Astra Ultra}
\date{9 October 2026}
\begin{document}
\maketitle
\noindent\textbf{Problem ID:} \texttt{D47-174}. \textbf{Status:} Unresolved; restricted results and reductions.

\section{Research target and the declared one-batch model}
The full target combines institutional withholding incentives, equity,
exchange failures, international constraints, and dynamic participation.
Roth's white paper \cite[p.~4]{roth} identifies full center submission as
an institutional problem; Ashlagi and Roth \cite[Section 6, item 2,
p.~5467]{ar} explicitly separate scale from biological heterogeneity.
Neither source is a theorem asserting the mechanism studied below.

We analyze a restricted, completely specified model. All pairs arrive
in one batch, the realized compatibility graph is known to the planner,
only cycles of length at most $K\geq2$ are permitted, and there are no
chains. Resources allow any collection of vertex-disjoint cycles;
there is no shared capacity constraint beyond vertex disjointness.
Each selected cycle succeeds in its entirety with the same probability
$s\in(0,1]$, and otherwise makes no transplant. Cycle failures may be
correlated, but their marginal probability is $s$ under both internal
and pooled operation. A plan is selected before failures and no repair
is attempted. These assumptions make expected transplant counts equal
to $s$ times planned counts. There are no monetary transfers.

Center $h$ owns $V_h$. Let $O_h$ be a maximum-cardinality internal cycle
packing and $b_h=|V(O_h)|$. Its outside option is $s b_h$. The union
$O=\bigcup_hO_h$ is feasible, and write $b=\sum_hb_h$. Computing these
packings may itself be hard for $K\geq3$; the guarantees below are
structural and do not assert a polynomial clearing algorithm.

\section{A guaranteed-participation benchmark}
Reserve the cycles $O$, delete their vertices, and choose a
maximum-cardinality packing $Q$ in the residual graph. Output $O\cup Q$.

\begin{theorem}[A cycle-cap bound on the cost of reservations]
For every realized graph in the stated model, the reservation plan gives
every center at least its internal expected outside option. If $T^*$ is
the maximum number of vertices covered by any pooled packing, then
\[
 |V(O\cup Q)|\geq\max\{b,T^*-(K-1)b\}\geq T^*/K.
\]
The same $1/K$ comparison holds for expected total transplants. If
patient weights satisfy $0<w_{\min}\leq w_i\leq w_{\max}$, the plan's
expected weighted value is at least
$w_{\min}/(K w_{\max})$ times the maximum expected weighted value.
\end{theorem}
\begin{proof}
Every reserved internal cycle is still performed with the same marginal
success probability, so center $h$ receives at least $s b_h$ expected
transplants. Take an optimum pooled packing $P^*$ with $T^*$ covered
vertices. At most $b$ of its cycles intersect the reserved vertices,
because the cycles are disjoint and each intersecting cycle contains
at least one distinct reserved vertex. Deleting these cycles discards
at most $K b$ vertices. The remaining cycles are a feasible residual
packing, so $|V(Q)|\geq T^*-K b$ and
$|V(O\cup Q)|\geq T^*-(K-1)b$. The trivial bound $|V(O\cup Q)|\geq b$
gives the maximum. If $b\geq T^*/K$ the first term suffices; otherwise
the second is greater than $T^*/K$.

Linearity of expectation multiplies all counts by $s$; independence
of failures is unnecessary. The algorithm's weighted value is at least
$s w_{\min}T^*/K$, while any packing's expected weighted value is at
most $s w_{\max}T^*$, proving the final assertion.
\end{proof}

This benchmark enforces participation conditional on full submission.
It is not a dominant-strategy withholding mechanism: changing the
submitted set can change the reserved internal optimum and the residual
exchange opportunities. The theorem deliberately makes no such
incentive claim.

\begin{proposition}[A truthful baseline and its efficiency limit]
The mechanism that independently performs a maximum-cardinality internal
packing for each submitting center is dominant for full submission
against subset withholding and subsequent internal exchange on withheld
vertices, with the same cycle cap, failure law, and resource rules.
Nevertheless its ratio to pooled expected total transplants can be zero.
\end{proposition}
\begin{proof}
Fix a center's full set $V_h$ and a submitted subset $S_h$. A maximum
internal packing on $S_h$ together with any packing on
$V_h\setminus S_h$ is a feasible internal packing on $V_h$. Their total
planned count cannot exceed $b_h$. Multiplication by $s$ proves the
required expected-utility dominance, for every submission of the
other centers; those submissions do not enter this mechanism.
For the efficiency assertion give each of two centers one vertex and
put a directed two-cycle between the vertices, with no internal cycles.
The mechanism makes zero transplants, whereas pooling gives $2s>0$.
\end{proof}

Thus a participation-safe allocation and a dominant-submission rule
are mathematically distinct. Combining their benefits, rather than
confusing the two inequalities, is part of the remaining target.

\section{An exact equity obstruction, including lotteries and failures}
\begin{proposition}[No uniform weighted efficiency under unrestricted weights]
Even with $K=2$ and the common success probability $s$, individually
rational randomized mechanisms can have an arbitrarily small ratio to
unconstrained weighted efficiency, irrespective of their computational
power or withholding incentives.
\end{proposition}
\begin{proof}
Center $A$ owns $a_1,a_2$ and center $B$ owns $b_1$. The only cycles are
$(a_1,a_2)$ and $(a_1,b_1)$. Set $w_{a_1}=w_{a_2}=1$ and $w_{b_1}=W>1$.
Center $A$ has outside option $2s$. A lottery putting probability $x$
on its internal cycle and $y$ on the cross-center cycle, with
$x+y\leq1$, gives $A$ expected utility $s(2x+y)$. The requirement
$2x+y\geq2$ forces $x=1,y=0$: indeed $2x+y\leq2-y$.
Thus every individually rational lottery has weighted value $2s$.
The unconstrained optimum uses the cross-center cycle and has value
$s(W+1)$. Their ratio $2/(W+1)$ tends to zero. All plans obey the same
failure and medical constraints, and no payment can compensate $A$.
\end{proof}

The loss is caused by the difference between center transplant-count
objectives and patient equity weights. Randomizing plans does not
remove the binding participation constraint. This is a restricted
impossibility result, not a prediction about actual patient weights.

\section{Separating composition and scale in a two-type pool}
To make the counterfactuals explicit, assume pairs have type $X$ or $Y$,
any opposite-type pair can form a two-cycle, and no same-type cycle is
feasible. Let a pool contain $x$ and $y$ vertices of the two types.
Failure probability remains $1-s$ per cycle.

\begin{theorem}[Exact composition accounting and a finite-pool bound]
The maximum planned transplant count is $f(x,y)=2\min(x,y)$.
For two pools with imbalances $d_j=x_j-y_j$, the exact gain from
pooling their fixed compositions is
\[
 f(x_1+x_2,y_1+y_2)-f(x_1,y_1)-f(x_2,y_2)
 =|d_1|+|d_2|-|d_1+d_2|.
\]
For an iid pool of size $N$ with type-$X$ probability $p$, let
$F(N,p)$ be its expected optimum planned count. Then
\[
 F(N,p)=N-\E|2X-N|,\qquad X\sim\operatorname{Bin}(N,p),
\]
\[
 \max\{0,N-N|2p-1|-2\sqrt{Np(1-p)}\}
 \leq F(N,p)\leq N-N|2p-1|.
\]
Multiply all formulas by $s$ for expected completed transplants.
\end{theorem}
\begin{proof}
Each cycle uses one vertex of each type, so at most $\min(x,y)$ cycles
are possible; completeness across types attains this bound. The identity
$2\min(x,y)=x+y-|x-y|$ proves the fixed-composition formula by subtraction.
For the random pool substitute $x=X,y=N-X$. Convexity of absolute value
gives $\E|2X-N|\geq|2Np-N|$. The triangle inequality and
Cauchy--Schwarz give
\[
 \E|2X-N|\leq |2Np-N|+2\E|X-Np|
 \leq N|2p-1|+2\sqrt{Np(1-p)}.
\]
Subtract from $N$ and use nonnegativity. Linearity of expectation again
supplies the common factor $s$.
\end{proof}

For a declared common composition $p_0$, a scale-only counterfactual is
$F(N_1+N_2,p_0)-F(N_1,p_0)-F(N_2,p_0)$. At the same pooled size, a
composition comparison is $F(N_1+N_2,p)-F(N_1+N_2,p_0)$.
A mixture of two countries with fixed counts and different probabilities
need not have binomial type counts; its corresponding expression uses
the sum of their two independent binomials. Choosing $p_0$ is a
substantive counterfactual, not an empirically identified decomposition.

\section{What remains}
These proofs handle batch arrivals, a transparent failure law, outside
options, and one explicit biological simplification. They neither
calibrate realistic blood/HLA composition nor solve dynamic exit,
heterogeneous failures, shared capacity, chains, or strategic pooling.
The unrestricted-weight obstruction persists even before those harder
features are added. Sustainable mechanisms with quantified withholding
incentives remain the principal missing step.
\begin{thebibliography}{9}
\bibitem{roth} A. E. Roth (2010). \emph{Market Design: Understanding
markets well enough to fix them when they're broken}, Section 1, p.~4.
\url{https://web.stanford.edu/~alroth/papers/NSF%20Grand%20Challenge.SBE%202020.%20Market%20Design.pdf}.
\bibitem{ar} I. Ashlagi and A. E. Roth (2021). \emph{Kidney Exchange:
An Operations Perspective}. Management Science 67(9), 5455--5478.
Sections 4.2 and 6, particularly item 2 on p.~5467.
\url{https://pubsonline.informs.org/doi/pdf/10.1287/mnsc.2020.3954}.
\end{thebibliography}

\end{document}
